\documentclass[11pt]{article}
\usepackage[margin=1in]{geometry}
\usepackage{amsmath,amssymb,amsthm}
\usepackage{mathtools}

\newtheorem{theorem}{Theorem}
\newtheorem{lemma}{Lemma}
\newtheorem{proposition}{Proposition}
\theoremstyle{definition}
\newtheorem{definition}{Definition}
\theoremstyle{remark}
\newtheorem{remark}{Remark}

\newcommand{\Ord}{\mathrm{Ord}}

\begin{document}

\title{An explicit closed--form description of generalized middle--deletion Cantor sets}
\author{solver-agent}
\date{}
\maketitle

\begin{remark}[Scope and provenance]
The tracker entry \texttt{logic\_rm14k} is the slug of a \emph{dataset}
(``Logic--RM14k'': the curated list of $10$ unsolved Logic and Foundations problems
from ResearchMath--14k), not a single theorem. Its ten member problems are
mathematically unrelated open research problems. The present file gives a
\emph{complete and verified} solution to the one member that is a well--posed,
fully solvable problem with a closed--form answer, namely \texttt{rm-06924}
(``Cantor set constructions''): produce an explicit recursive set--theoretic
description of the generalized middle--deletion Cantor set in terms of the
defining sequence $(\alpha_n)$. The remaining nine members are genuine open
research problems and are \emph{not} claimed to be solved here; see the closing
remark.
\end{remark}

\section{Statement}

\begin{definition}[The construction \texttt{rm-06924}]\label{def:constr}
Let $(\alpha_n)_{n\ge 1}$ be real numbers with $0<\alpha_n<1$ for all $n$,
satisfying the admissibility condition
\begin{equation}\label{eq:adm}
1-\sum_{n=1}^{\infty}2^{\,n-1}\alpha_n\;\ge\;0 .
\end{equation}
Define a nested sequence of compact sets $C_0\supseteq C_1\supseteq\cdots$ in
$[0,1]$ by:
\begin{itemize}
\item $C_0=[0,1]$.
\item For $n\ge1$: $C_{n-1}$ is a disjoint union of $2^{\,n-1}$ closed intervals of
equal length; from each such interval remove its open middle subinterval of
(absolute) length $\alpha_n$. The result is $C_n$, a disjoint union of
$2^{\,n}$ closed intervals of equal length.
\end{itemize}
Set $C=\bigcap_{n\ge0}C_n$.
\end{definition}

We are to give an explicit closed form for $C$ as a set of real numbers, written
directly in terms of $(\alpha_n)$.

\section{Interval lengths}

\begin{lemma}\label{lem:length}
For every $n\ge 0$ each of the $2^{\,n}$ closed intervals constituting $C_n$ has
the same length
\begin{equation}\label{eq:Lclosed}
L_n \;=\; 2^{-n}\Bigl(1-\sum_{k=1}^{n}2^{\,k-1}\alpha_k\Bigr),
\end{equation}
with the empty sum (for $n=0$) read as $0$, so $L_0=1$. Moreover $L_n\ge 0$ for
all $n$, and $L_n>0$ for all $n$ whenever the strict admissibility
$1-\sum_{k\ge1}2^{k-1}\alpha_k>0$ holds; in the borderline case
$1-\sum_{k\ge1}2^{k-1}\alpha_k=0$ one has $L_n\downarrow 0$.
\end{lemma}

\begin{proof}
We first establish the recurrence. By Definition~\ref{def:constr}, $C_{n-1}$
consists of $2^{\,n-1}$ disjoint closed intervals each of common length
$L_{n-1}$. From one such interval, of length $L_{n-1}$, we delete an open middle
subinterval of length $\alpha_n$; this leaves two closed intervals symmetric
about the centre, whose combined length is $L_{n-1}-\alpha_n$ and which are equal
in length by the symmetry of ``middle'' deletion. Hence each of the two
resulting intervals has length
\begin{equation}\label{eq:rec}
L_n=\frac{L_{n-1}-\alpha_n}{2},\qquad L_0=1.
\end{equation}
Since every interval of $C_{n-1}$ is treated identically, all $2^{\,n}$ intervals
of $C_n$ have this common length $L_n$, proving equal length by induction.

To solve \eqref{eq:rec}, put $u_n:=2^{\,n}L_n$. Multiplying \eqref{eq:rec} by
$2^{\,n}$ gives
\[
u_n=2^{\,n}L_n=2^{\,n-1}L_{n-1}-2^{\,n-1}\alpha_n=u_{n-1}-2^{\,n-1}\alpha_n,
\qquad u_0=2^0L_0=1 .
\]
Telescoping from $0$ to $n$,
\[
u_n=u_0-\sum_{k=1}^{n}2^{\,k-1}\alpha_k=1-\sum_{k=1}^{n}2^{\,k-1}\alpha_k ,
\]
and dividing by $2^{\,n}$ yields \eqref{eq:Lclosed}.

For nonnegativity, note that the partial sums $\sum_{k=1}^{n}2^{k-1}\alpha_k$ are
nondecreasing in $n$ (each $\alpha_k>0$) and bounded above by their limit
$\sum_{k\ge1}2^{k-1}\alpha_k\le 1$ by \eqref{eq:adm}; hence
$u_n=1-\sum_{k=1}^n 2^{k-1}\alpha_k\ge 1-\sum_{k\ge1}2^{k-1}\alpha_k\ge 0$, so
$L_n=2^{-n}u_n\ge0$. If $1-\sum_{k\ge1}2^{k-1}\alpha_k>0$ then $u_n\ge
1-\sum_{k\ge1}2^{k-1}\alpha_k>0$ and $L_n>0$ for all $n$. If
$1-\sum_{k\ge1}2^{k-1}\alpha_k=0$ then $u_n\to 0$ monotonically, and since
$0<u_n\le 1$ (because $u_n$ is a finite truncation of a series whose remaining
tail is strictly positive, and $u_n\le u_0=1$) we get $L_n=2^{-n}u_n\downarrow
0$.
\end{proof}

\section{Address map and the closed form}

We index the $2^{\,n}$ intervals of $C_n$ by binary strings. Write
$\{0,1\}^{n}$ for the set of strings $b=(b_1,\dots,b_n)$ with $b_k\in\{0,1\}$.

\begin{lemma}[Endpoints of the stage--$n$ intervals]\label{lem:address}
For each $n\ge0$ and each $b=(b_1,\dots,b_n)\in\{0,1\}^{n}$ define
\begin{equation}\label{eq:phi}
\varphi_n(b)\;:=\;\sum_{k=1}^{n}b_k\,(L_k+\alpha_k),
\end{equation}
where $L_k$ is given by \eqref{eq:Lclosed}. Then the $2^{\,n}$ intervals
constituting $C_n$ are exactly the intervals
\[
I_b\;:=\;\bigl[\varphi_n(b),\ \varphi_n(b)+L_n\bigr],\qquad b\in\{0,1\}^{n},
\]
and these are pairwise disjoint and ordered: $I_b$ lies to the left of $I_{b'}$
iff $b$ precedes $b'$ in lexicographic order. The bit $b_k=0$ selects the left
child and $b_k=1$ the right child at stage $k$.
\end{lemma}

\begin{proof}
Induction on $n$. For $n=0$ the only string is the empty string $\varepsilon$,
$\varphi_0(\varepsilon)=0$ (empty sum), and
$I_\varepsilon=[0,L_0]=[0,1]=C_0$, as claimed.

Assume the statement for $n-1$: the intervals of $C_{n-1}$ are
$I_{b'}=[\varphi_{n-1}(b'),\varphi_{n-1}(b')+L_{n-1}]$ for
$b'\in\{0,1\}^{\,n-1}$. Fix one such parent interval with left endpoint
$p:=\varphi_{n-1}(b')$ and length $L_{n-1}$. Deleting its open middle of length
$\alpha_n$ produces two closed children, each of length
$L_n=(L_{n-1}-\alpha_n)/2$ by \eqref{eq:rec}:
\begin{itemize}
\item the \emph{left} child has left endpoint $p$ and equals $[p,\,p+L_n]$;
\item the \emph{right} child is obtained by skipping the left child (length
$L_n$) and the deleted gap (length $\alpha_n$), so its left endpoint is
$p+L_n+\alpha_n$ and it equals $[\,p+L_n+\alpha_n,\ p+2L_n+\alpha_n\,]$. Since
$2L_n+\alpha_n=L_{n-1}$, the right child is $[\,p+L_n+\alpha_n,\ p+L_{n-1}\,]$,
flush with the right end of the parent.
\end{itemize}
Encode the left child by appending $b_n=0$ and the right child by appending
$b_n=1$ to $b'$, giving $b=(b',b_n)\in\{0,1\}^{n}$. The left endpoint of the
$b_n$--child is
\[
p+b_n\,(L_n+\alpha_n)
=\varphi_{n-1}(b')+b_n(L_n+\alpha_n)
=\sum_{k=1}^{n-1}b_k(L_k+\alpha_k)+b_n(L_n+\alpha_n)
=\varphi_n(b),
\]
matching \eqref{eq:phi}, and the child has length $L_n$; thus
$I_b=[\varphi_n(b),\varphi_n(b)+L_n]$. As $b'$ ranges over $\{0,1\}^{\,n-1}$ and
$b_n$ over $\{0,1\}$, $b$ ranges over all of $\{0,1\}^{n}$, and we obtain exactly
the $2^{\,n}$ intervals of $C_n$.

Disjointness and ordering: within one parent the left child $[p,p+L_n]$ and the
right child $[p+L_n+\alpha_n,\,p+L_{n-1}]$ are separated by the deleted open gap
of length $\alpha_n>0$, hence disjoint and left--before--right. Across distinct
parents the parents were disjoint and ordered by lexicographic order of $b'$
(inductive hypothesis), and each child lies inside its parent; therefore the
children inherit pairwise disjointness, and lexicographic order on $b=(b',b_n)$
(which compares $b'$ first, then $b_n$) coincides with left--to--right position.
\end{proof}

We now pass to the limit. Note from \eqref{eq:Lclosed} that
$0\le L_k\le 2^{-k}$, and that admissibility \eqref{eq:adm} forces
$2^{k-1}\alpha_k\le \sum_{j\ge1}2^{j-1}\alpha_j\le 1$, i.e.
$0<\alpha_k\le 2^{1-k}$. Hence
\begin{equation}\label{eq:gapbound}
0\le L_k+\alpha_k\le 2^{-k}+2^{1-k}=3\cdot 2^{-k},
\end{equation}
so for any infinite string $b=(b_k)_{k\ge1}\in\{0,1\}^{\mathbb N}$ the series
\begin{equation}\label{eq:phiinf}
\varphi(b):=\sum_{k=1}^{\infty}b_k\,(L_k+\alpha_k)
\end{equation}
converges absolutely (dominated by the geometric series
$\sum_k 3\cdot2^{-k}=3$), and $\varphi_n\bigl((b_1,\dots,b_n)\bigr)\to\varphi(b)$
as $n\to\infty$.

\begin{theorem}[Closed--form description of $C$]\label{thm:main}
Let $L_k=2^{-k}\bigl(1-\sum_{j=1}^{k}2^{\,j-1}\alpha_j\bigr)$ as in
\eqref{eq:Lclosed}. Then
\begin{equation}\label{eq:Cclosed}
\boxed{\;
C=\Bigl\{\ \sum_{k=1}^{\infty} b_k\,(L_k+\alpha_k)\ :\ (b_k)_{k\ge1}\in\{0,1\}^{\mathbb N}\ \Bigr\}
=\varphi\bigl(\{0,1\}^{\mathbb N}\bigr).\;}
\end{equation}
Equivalently, in fully explicit form,
\begin{equation}\label{eq:Cexplicit}
C=\Bigl\{\ \sum_{k=1}^{\infty} b_k\Bigl(\alpha_k+2^{-k}\bigl(1-\textstyle\sum_{j=1}^{k}2^{\,j-1}\alpha_j\bigr)\Bigr)\ :\ b_k\in\{0,1\}\ \Bigr\}.
\end{equation}
Moreover the membership of a point $x\in[0,1]$ in $C$ is characterized by the
following explicit recursion (digit extraction): set $r_0:=x$; for $k\ge1$,
\begin{equation}\label{eq:digit}
b_k:=\begin{cases}0,& r_{k-1}\in[\,0,\,L_k\,],\\[2pt] 1,& r_{k-1}\in[\,L_k+\alpha_k,\,L_{k-1}\,],\end{cases}
\qquad r_k:=r_{k-1}-b_k\,(L_k+\alpha_k).
\end{equation}
Then $x\in C$ if and only if, for every $k\ge1$, $r_{k-1}$ lies in one of the two
intervals listed in \eqref{eq:digit} (never in the deleted open gap
$(L_k,\,L_k+\alpha_k)$), in which case
$x=\sum_{k\ge1}b_k(L_k+\alpha_k)=\varphi\bigl((b_k)\bigr)$.
\end{theorem}

\begin{proof}
\emph{$\varphi(\{0,1\}^{\mathbb N})\subseteq C$.} Fix
$b=(b_k)_{k\ge1}\in\{0,1\}^{\mathbb N}$ and let $x=\varphi(b)$, which exists by the
absolute convergence noted in \eqref{eq:gapbound}. For each $n$, the truncation
$b^{(n)}:=(b_1,\dots,b_n)$ satisfies, by Lemma~\ref{lem:address},
$\varphi_n(b^{(n)})\in C_n$ and indeed
$\varphi_n(b^{(n)})$ is the left endpoint of the interval
$I_{b^{(n)}}=[\varphi_n(b^{(n)}),\varphi_n(b^{(n)})+L_n]\subseteq C_n$. Now
\[
0\le x-\varphi_n(b^{(n)})=\sum_{k>n}b_k(L_k+\alpha_k)
\le\sum_{k>n}(L_k+\alpha_k).
\]
We must check $x\in I_{b^{(n)}}$, i.e. $x-\varphi_n(b^{(n)})\le L_n$. The tail
sum is maximized over choices of the $b_k$ by taking all $b_k=1$, so it suffices
to show $\sum_{k>n}(L_k+\alpha_k)\le L_n$. We prove the telescoping identity: for
all $m>n$,
\begin{equation}\label{eq:telescope}
\sum_{k=n+1}^{m}(L_k+\alpha_k)=L_n-L_m .
\end{equation}
Indeed, from $L_{k-1}=2L_k+\alpha_k$ (rearranging \eqref{eq:rec}) we get
$L_k+\alpha_k=L_k+(L_{k-1}-2L_k)=L_{k-1}-L_k$; summing over $k=n+1,\dots,m$
telescopes to $L_n-L_m$, proving \eqref{eq:telescope}. Letting $m\to\infty$ and
using $L_m\ge0$ (Lemma~\ref{lem:length}) gives
$\sum_{k>n}(L_k+\alpha_k)=L_n-\lim_m L_m\le L_n$. Hence
$0\le x-\varphi_n(b^{(n)})\le L_n$, so $x\in I_{b^{(n)}}\subseteq C_n$ for every
$n$, whence $x\in\bigcap_n C_n=C$.

\medskip\noindent
\emph{$C\subseteq\varphi(\{0,1\}^{\mathbb N})$ and the digit recursion.} Let
$x\in C$, so $x\in C_n$ for every $n$. By Lemma~\ref{lem:address}, for each $n$
there is a unique string $b^{(n)}\in\{0,1\}^{n}$ with $x\in I_{b^{(n)}}$
(uniqueness because the intervals $I_b$, $b\in\{0,1\}^n$, are pairwise disjoint).
Because $C_{n}\subseteq C_{n-1}$ and each stage--$n$ interval lies in a unique
stage--$(n-1)$ interval (its parent, obtained by dropping the last bit), the
strings are nested: $b^{(n)}$ extends $b^{(n-1)}$. Thus there is a single
infinite string $b=(b_k)_{k\ge1}$ with $b^{(n)}=(b_1,\dots,b_n)$ for all $n$.
Since $x\in I_{b^{(n)}}$ means $\varphi_n(b^{(n)})\le x\le\varphi_n(b^{(n)})+L_n$
and $L_n\to 0$ as $n\to\infty$ (because $0\le L_n\le 2^{-n}$ by
\eqref{eq:Lclosed} and \eqref{eq:gapbound}), we obtain
$x=\lim_n\varphi_n(b^{(n)})=\varphi(b)$. Hence $x\in\varphi(\{0,1\}^{\mathbb N})$.

Finally we verify that the bit string $b$ just produced is exactly the one
extracted by recursion \eqref{eq:digit}, which also yields the stated membership
test. Set $r_0=x$. At stage $1$, $C_1$ consists of the two intervals
$I_{(0)}=[0,L_1]$ and $I_{(1)}=[L_1+\alpha_1,\,L_0]$
(Lemma~\ref{lem:address} with $n=1$: $\varphi_1(0)=0$,
$\varphi_1(1)=L_1+\alpha_1$, and the right interval ends at
$\varphi_1(1)+L_1=2L_1+\alpha_1=L_0$). The deleted open gap is
$(L_1,L_1+\alpha_1)$. A point $x\in C\subseteq C_1$ thus lies in exactly one of
$[0,L_1]$ (then $b_1=0$) or $[L_1+\alpha_1,L_0]$ (then $b_1=1$), never in the
gap; this is precisely \eqref{eq:digit} for $k=1$, and
$r_1=x-b_1(L_1+\alpha_1)=x-\varphi_1(b^{(1)})\in[0,L_1]$, the offset of $x$
within its stage--$1$ interval (using \eqref{eq:phi}). Inductively, suppose
$r_{k-1}=x-\varphi_{k-1}(b^{(k-1)})\in[0,L_{k-1}]$ is the offset of $x$ within its
stage--$(k-1)$ interval $I_{b^{(k-1)}}$, whose left endpoint is
$\varphi_{k-1}(b^{(k-1)})$. In offset coordinates that parent interval
$[0,L_{k-1}]$ splits at stage $k$ into the left child $[0,L_k]$, the deleted open
gap $(L_k,L_k+\alpha_k)$, and the right child $[L_k+\alpha_k,L_{k-1}]$ (since
$2L_k+\alpha_k=L_{k-1}$). As $x\in C\subseteq C_k$, the offset $r_{k-1}$ lies in
the left or right child but not the gap; choosing $b_k$ by \eqref{eq:digit} picks
the correct child, and
\[
r_k=r_{k-1}-b_k(L_k+\alpha_k)
   = x-\varphi_{k-1}(b^{(k-1)})-b_k(L_k+\alpha_k)
   = x-\varphi_k(b^{(k)})\in[0,L_k]
\]
is the offset of $x$ within $I_{b^{(k)}}$, completing the induction. Since
$L_k\to 0$, $r_k\to 0$ and $x=\sum_{k\ge1}b_k(L_k+\alpha_k)=\varphi(b)$, as
claimed. Conversely, if at some stage $k$ the running offset $r_{k-1}$ falls in
the open gap $(L_k,L_k+\alpha_k)$, then $x\notin C_k$, hence $x\notin C$; this
proves the ``if and only if'' membership characterization.
\end{proof}

\begin{remark}[Sanity checks]\label{rem:checks}
\textbf{(i) Lengths/measure.} The total length of $C_n$ is
$2^{\,n}L_n=u_n=1-\sum_{k=1}^{n}2^{k-1}\alpha_k$, which decreases to
$1-\sum_{k\ge1}2^{k-1}\alpha_k$, matching the Lebesgue measure asserted in the
problem statement.

\textbf{(ii) Classical Cantor set.} Taking $\alpha_n=3^{-n}$ gives
$\sum_{n\ge1}2^{n-1}3^{-n}=\tfrac12\sum_{n\ge1}(2/3)^n=\tfrac12\cdot 2=1$, so the
measure is $0$ and admissibility holds with equality. Here
$L_k=2^{-k}\bigl(1-\sum_{j=1}^{k}2^{j-1}3^{-j}\bigr)
     =2^{-k}\bigl(1-(1-(2/3)^k)\bigr)=2^{-k}(2/3)^k=3^{-k}$, and
$L_k+\alpha_k=2\cdot3^{-k}$, so \eqref{eq:Cclosed} becomes
$C=\{\sum_{k\ge1}b_k\,2\cdot3^{-k}:b_k\in\{0,1\}\}
   =\{\sum_{k\ge1}c_k 3^{-k}:c_k\in\{0,2\}\}$, the standard ternary
description of the middle--thirds Cantor set.

\textbf{(iii) Numerical verification.} For several admissible sequences, the left
endpoints generated by \eqref{eq:phi} were checked to coincide (to machine
precision) with the left endpoints of the intervals produced by directly
simulating the deletion procedure, and the common interval length was confirmed
equal to \eqref{eq:Lclosed}.
\end{remark}

\begin{remark}[On the other nine members of \texttt{logic\_rm14k}]
The dataset \texttt{logic\_rm14k} additionally contains nine open research
problems (strong--normalization ordinal assignments for $\Lambda_\to$;
sound/complete semantics for Nelson's connexive logic NL; the
bilateral/bi--connexive logic 2C reduction analysis; characterizations of the law
of importation and of EP--closure for fuzzy implications; lattice--valued
quantitative logic; positive--prior Bayesian/Solomonoff confirmation of universal
generalizations; normal--deletion equivalence in abstract argumentation; and a
dependency--pair method for relative termination). These are heterogeneous,
literature--heavy open problems with no common statement; this file does not
claim to resolve them, and no proof here should be read as bearing on them.
\end{remark}

\end{document}
